\documentclass[12pt,a4paper]{article}

% ==================================================
% 1. CONFIGURATION
% ==================================================

% Enable / disable title page
\newif\ifshowtitlepage
\showtitlepagetrue
% \showtitlepagefalse

% Enable / disable table of contents
\newif\ifshowtoc
\showtoctrue

% Enable / disable bibliography
\newif\ifshowbibliography
\showbibliographyfalse

% Document information
\newcommand{\worktitle}{Definition of Done}
\newcommand{\worksubtitle}{NAO Robot Exercise System}
\newcommand{\teamname}{NAOMI}
\newcommand{\subject}{NAO Robot Exercise System}

\newcommand{\university}{Slovak University of Technology in Bratislava}
\newcommand{\faculty}{Faculty of Informatics and Information Technologies}

\newcommand{\academicyear}{2026/2027}
\newcommand{\location}{Bratislava}
\newcommand{\documentyear}{2026/2027}

% Team leaders
\newcommand{\leaderone}{Ing. Jakub Perdek, PhD.}
\newcommand{\leadertwo}{ doc. Ing. Fedor Lehocki PhD.}

% Team members
\newcommand{\memberone}{Bc. Akhmadinurov Amal}
\newcommand{\membertwo}{Bc. Danilov Denis}
\newcommand{\memberthree}{Bc. Dmitriev Anton}
\newcommand{\memberfour}{Bc. Gainetdinov Ruslan}
\newcommand{\memberfive}{Bc. Kuzin Timofei}
\newcommand{\membersix}{Bc. Tikhoblazhenko Alena}
\newcommand{\memberseven}{Bc. Zimnitski Aliaksei}


% ==================================================
% 2. PACKAGES
% ==================================================

\usepackage[
    a4paper,
    top=3cm,
    bottom=2.7cm,
    left=2.8cm,
    right=2.8cm,
    headheight=48pt,
    headsep=20pt
]{geometry}

\usepackage[T1]{fontenc}
\usepackage[utf8]{inputenc}
\usepackage[english]{babel}
\usepackage{lmodern}

\usepackage{graphicx}
\usepackage[table]{xcolor}
\usepackage{fancyhdr}
\usepackage{titlesec}
\usepackage{setspace}
\usepackage{lastpage}
\usepackage{array}
\usepackage{amssymb}
\usepackage{float}
\usepackage{caption}
\usepackage{booktabs}
\usepackage{tabularx}
\usepackage{lipsum}
\usepackage{csquotes}

% Bibliography
\usepackage[
    backend=biber,
    style=ieee,
    sorting=none
]{biblatex}

\addbibresource{references.bib}

% Hyperlinks
\usepackage[hidelinks]{hyperref}


% ==================================================
% 3. COLORS
% ==================================================

\definecolor{primary}{HTML}{000000}
\definecolor{secondary}{HTML}{000000}
\definecolor{muted}{HTML}{000000}
\definecolor{linecolor}{HTML}{000000}
\definecolor{naomiTableHeader}{HTML}{17445A}
\definecolor{naomiTableStripe}{HTML}{EDF7FA}
\definecolor{naomiTableRule}{HTML}{66BBD0}

% Ragged-right columns keep narrow table cells readable.
\newcolumntype{L}[1]{>{\raggedright\arraybackslash}p{#1}}
\newcolumntype{Y}{>{\raggedright\arraybackslash}X}


% ==================================================
% 4. TYPOGRAPHY
% ==================================================

\setstretch{1.2}

\setlength{\parindent}{0pt}
\setlength{\parskip}{7pt}

% Section
\titleformat{\section}
    {\Large\bfseries\color{primary}}
    {\thesection}
    {0.8em}
    {}

% Subsection
\titleformat{\subsection}
    {\large\bfseries\color{primary}}
    {\thesubsection}
    {0.8em}
    {}

% Subsubsection
\titleformat{\subsubsection}
    {\normalsize\bfseries\color{primary}}
    {\thesubsubsection}
    {0.8em}
    {}

\titlespacing*{\section}{0pt}{20pt}{10pt}
\titlespacing*{\subsection}{0pt}{14pt}{7pt}
\titlespacing*{\subsubsection}{0pt}{12pt}{5pt}

% Captions
\captionsetup{
    font=small,
    labelfont=bf,
    justification=centering
}

% Bibliography title
\DefineBibliographyStrings{english}{
    bibliography = {References}
}


% ==================================================
% 5. HEADER AND FOOTER
% ==================================================

\pagestyle{fancy}
\fancyhf{}

% Header logo
\fancyhead[L]{
    \includegraphics[
        width=4cm,
        height=1.15cm,
        keepaspectratio
    ]{assets/horizontal.png}
}

% Header information
\fancyhead[R]{
    \begin{tabular}{r}
        {\footnotesize\bfseries \worksubtitle}\\[3pt]
        {\scriptsize Academic Year \academicyear}
    \end{tabular}
}

% Header line
\renewcommand{\headrulewidth}{0.7pt}

% Footer left
\fancyfoot[L]{
    {\scriptsize
    \teamname\enspace|\enspace\subject}
}

% Footer right
\fancyfoot[R]{
    {\scriptsize
    \thepage\enspace/\enspace\pageref{LastPage}}
}

% Footer line
\renewcommand{\footrulewidth}{0.4pt}

% Plain pages
\fancypagestyle{plain}{
    \fancyhf{}

    \fancyhead[L]{
        \includegraphics[
            width=4cm,
            height=1.15cm,
            keepaspectratio
        ]{assets/horizontal.png}
    }

    \fancyhead[R]{
        \begin{tabular}{r}
            {\footnotesize\bfseries \worksubtitle}\\[3pt]
            {\scriptsize Academic Year \academicyear}
        \end{tabular}
    }

    \fancyfoot[L]{
        {\scriptsize
        \teamname\enspace|\enspace\subject}
    }

    \fancyfoot[R]{
        {\scriptsize
        \thepage\enspace/\enspace\pageref{LastPage}}
    }

    \renewcommand{\headrulewidth}{0.7pt}
    \renewcommand{\footrulewidth}{0.4pt}
}




% ==================================================
% 6. TITLE PAGE
% ==================================================

\newcommand{\makemytitlepage}{
    \begin{titlepage}
        \thispagestyle{empty}
        \centering

        % ==========================================
        % UNIVERSITY LOGO
        % ==========================================

        \vspace*{0.3cm}

        \includegraphics[
            width=6cm,
            height=5cm,
            keepaspectratio
        ]{assets/vertical.png}

        \vspace{0.5cm}

        {\large\bfseries
        \university\par}

        \vspace{0.2cm}

        {\normalsize
        \faculty\par}

        % ==========================================
        % PROJECT TITLE
        % ==========================================

        \vspace{1.7cm}

        {\small\MakeUppercase{\worksubtitle}\par}

        \vspace{0.5cm}

        \begin{minipage}{0.95\textwidth}
            \centering

            {\fontsize{25}{32}\selectfont\bfseries
            \worktitle\par}
        \end{minipage}

        % ==========================================
        % TEAM INFORMATION
        % ==========================================

        \vfill

        \begin{minipage}{0.9\textwidth}
            \centering
            \begin{minipage}[t]{0.46\linewidth}
                \centering
                {\normalsize\bfseries Team Leaders\par}
                \vspace{0.25cm}
                {\small
                \leaderone\par
                \leadertwo\par}
            \end{minipage}%
            \hfill
            \begin{minipage}[t]{0.48\linewidth}
                \centering
                {\normalsize\bfseries Team Members\par}
                \vspace{0.25cm}
                {\small
                \memberone\par
                \membertwo\par
                \memberthree\par
                \memberfour\par
                \memberfive\par
                \membersix\par
                \memberseven\par}
            \end{minipage}
        \end{minipage}

        % ==========================================
        % BOTTOM INFORMATION
        % ==========================================

        \vfill

        {\small
        \location\enspace\textbullet\enspace\academicyear\par}

    \end{titlepage}
}

% ==================================================
% 7. IMAGE PLACEHOLDER
% ==================================================

% Usage:
% \imageplaceholder{Text}{Caption}

\newcommand{\imageplaceholder}[2]{

    \begin{figure}[H]
        \centering

        \begingroup
        \setlength{\fboxsep}{0pt}

        \fbox{
            \begin{minipage}[c][5cm][c]{0.85\textwidth}

                \centering

                {\Large\bfseries IMAGE}

                \vspace{0.3cm}

                {\normalsize #1}

            \end{minipage}
        }

        \endgroup

        \caption{#2}

    \end{figure}
}


% ==================================================
% 8. REAL IMAGE COMMAND
% ==================================================

% Usage:
% \documentimage{path}{caption}{label}

\newcommand{\documentimage}[3]{

    \begin{figure}[htbp]
        \centering

        \includegraphics[
            width=0.85\textwidth,
            keepaspectratio
        ]{#1}

        \caption{#2}
        \label{#3}

    \end{figure}
}


% ==================================================
% 9. TABLES
% ==================================================

% Usage:
% \begin{naomitable}{Caption}{label}
%     \begin{tabularx}{\textwidth}{L{2.8cm} Y L{2.6cm}}
%         \rowcolor{naomiTableHeader}
%         \textcolor{white}{\bfseries Name} & ... \\
%         ...
%         \arrayrulecolor{naomiTableRule}\bottomrule[1pt]
%     \end{tabularx}
% \end{naomitable}

\newenvironment{naomitable}[2]{%
    \begin{table}[htbp]
        \centering
        \caption{#1}
        \label{#2}
        \begingroup
        \small
        \renewcommand{\arraystretch}{1.4}
        \setlength{\tabcolsep}{7pt}
        \rowcolors{2}{naomiTableStripe}{white}
}{%
        \endgroup
    \end{table}
}

% Usage:
% \tableplaceholder{Caption}{label}

\newcommand{\tableplaceholder}[2]{
    \begin{naomitable}{#1}{#2}
        \begin{tabularx}{\textwidth}{L{2.8cm} Y L{2.6cm}}
            \rowcolor{naomiTableHeader}
            \textcolor{white}{\bfseries Name} &
            \textcolor{white}{\bfseries Description} &
            \textcolor{white}{\bfseries Status} \\
            Item A & Sample description & Completed \\
            Item B & Sample description & In Progress \\
            Item C & Sample description & Planned \\
            \arrayrulecolor{naomiTableRule}\bottomrule[1pt]
        \end{tabularx}
    \end{naomitable}
}


% ==================================================
% 10. DOCUMENT
% ==================================================

\begin{document}

\ifshowtitlepage
    \makemytitlepage
\fi

\ifshowtoc
    \tableofcontents
    \clearpage
\fi

% ==================================================
% 1. PURPOSE
% ==================================================

\section{Purpose of the Document}

This document defines the \textbf{Definition of Done (DoD)} for the NAOMI team.
It establishes a common understanding of when work can be considered complete
within the NAO Robot Exercise System project. The DoD helps the team maintain
consistent quality, track progress in Jira, and avoid marking unfinished work as
completed.

NAOMI follows a Scrumban approach. Issues are managed on a Jira board using
five states: \textbf{To Do, In Progress, In Review, Testing, and Done}.
The DoD describes the conditions that must be met before an issue is moved
to the final state.

% ==================================================
% 2. SCOPE
% ==================================================

\section{Scope and Responsibilities}

The DoD applies to development tasks, bug fixes, documentation, research,
experiments, and other deliverables tracked in Jira. Not every technical
criterion applies to every task: for example, a documentation-only task does
not require new source code. However, every completed task must meet its
acceptance criteria and provide a verifiable result.

The \textbf{assignee} is responsible for completing the work, keeping Jira
up to date, and providing evidence of completion. The \textbf{reviewers}
check the result independently. The team confirms that the necessary reviews
and tests have been completed before closing the issue.

% ==================================================
% 3. GENERAL CRITERIA
% ==================================================

\section{General Definition of Done}

A Jira issue can be moved to \textbf{Done} only when all applicable criteria
below have been satisfied:

\begin{itemize}
    \item The intended result is complete and meets the issue's acceptance
          criteria.
    \item The output can be reviewed or verified by another team member.
    \item Relevant changes, decisions, and results are recorded in Jira or
          linked documentation.
    \item Required reviews have been completed and important comments
          have been resolved.
    \item Appropriate testing or validation has been performed, with no
          known unresolved critical defects caused by the change.
    \item Documentation has been created or updated when necessary.
    \item The Jira issue contains a link to the final result, such as a
          pull request, document, experiment report, or demonstration.
    \item The issue status accurately reflects the actual progress.
\end{itemize}

If a requirement cannot be satisfied, the issue must not be marked Done
without an explicit team decision and a documented follow-up.
% ==================================================
% 4. DEVELOPMENT WORK
% ==================================================

\section{Definition of Done for Development Tasks}

For tasks involving source code, the following criteria apply:

\begin{itemize}
    \item The implementation meets the agreed requirements
          and follows the project's coding conventions.
    \item The code is committed to a GitHub branch linked
          to the corresponding Jira issue.
    \item Necessary tests have been completed successfully.
    \item Relevant documentation has been updated
          where required.
    \item All required automated checks pass, if available.
    \item The changes have been merged into the appropriate
          target branch following the review process
          described below.
\end{itemize}

\subsection{Pull Request Review and Approval}

All code changes must be submitted through a Pull Request
before being merged into the main branch.

The approval requirements depend on the PR author:

\begin{itemize}
    \item \textbf{Repository Owners and Maintainers:}
          Members responsible for the repository may merge
          their own Pull Requests without additional approvals.

    \item \textbf{Other Team Members:}
          Pull Requests require two approvals:
          one from a repository owner or maintainer
          and one from another assigned reviewer.
\end{itemize}

Before merging, all requested changes must be addressed
and blocking review comments resolved. Reviewers should
verify correctness, code quality, and compliance with
project requirements.

Repository owners and maintainers remain responsible
for the quality of changes they merge without
additional review.

% ==================================================
% 5. REVIEW AND TESTING
% ==================================================

\section{Review and Testing Process}

The \textbf{In Review} and \textbf{Testing} states represent two different
checks. Review focuses on implementation quality and completeness; testing
focuses on whether the result behaves as expected.

\begin{naomitable}{Verification stages before completion}{tab:dod-stages}
    \begin{tabularx}{\textwidth}{L{2.4cm} Y Y}
        \rowcolor{naomiTableHeader}
        \textcolor{white}{\bfseries Jira status} &
        \textcolor{white}{\bfseries Main activity} &
        \textcolor{white}{\bfseries Exit condition} \\
        In Review &
        Review the implementation or deliverable and provide feedback. &
        Required reviews are completed and blocking comments resolved. \\
        Testing &
        Execute relevant functional, integration, or manual checks. &
        Expected behaviour is verified and significant defects are resolved. \\
        Done &
        Confirm the issue meets the applicable DoD criteria. &
        Evidence is linked and the final result is ready to use. \\
        \arrayrulecolor{naomiTableRule}\bottomrule[1pt]
    \end{tabularx}
\end{naomitable}

For the NAO robot project, testing may include verifying motion sequences,
robot behaviour, input processing, or communication between components,
depending on the task. 

If testing reveals a defect, the issue should return to an appropriate
working state or a related bug issue should be created. Moving an issue
to Done does not replace the need to record unresolved problems.

% ==================================================
% 6. NON-CODE WORK
% ==================================================

\section{Definition of Done for Non-Code Tasks}

\subsection{Documentation}

A documentation task is complete when its content is accurate, readable,
consistent with the current implementation, and stored in the agreed
team location. The relevant team member has reviewed it, and any essential
corrections have been applied. Meeting notes are considered complete once
the main decisions, action items, responsible people, and relevant Jira
links have been recorded and shared with the team.

\subsection{Research and Experiments}

Research or experiment tasks are complete when the objective, method,
results, and conclusions are documented. The result must be accessible to
other team members, including any limitations or remaining uncertainties.
A negative or inconclusive result can still be Done if the agreed work was
carried out and its outcome is documented.

% ==================================================
% 7. JIRA CHECKLIST
% ==================================================

\section{Completion Checklist in Jira}

Before moving a task to Done, the assignee should verify the following
short checklist. Items that do not apply should be marked as not applicable
rather than silently ignored.

\begin{itemize}
    \item[$\square$] Acceptance criteria are met.
    \item[$\square$] The result is linked in Jira.
    \item[$\square$] Review and required approvals are complete.
    \item[$\square$] Relevant tests or validation have passed.
    \item[$\square$] Documentation is up to date.
    \item[$\square$] There are no unresolved blocking problems.
    \item[$\square$] The final status in Jira is correct.
\end{itemize}

\section{Maintaining the Definition of Done}

The team may update this definition during retrospectives or when the
project's tools and workflow change. Any change should be discussed with
the team, documented, and applied consistently to future work. The DoD
should remain practical: its purpose is to support reliable delivery, not
to create unnecessary administrative steps.

\ifshowbibliography
    \clearpage
    \printbibliography[heading=bibintoc,title={References}]
\fi

\end{document}
